\heading{数学物理方法（下）-课程说明}
\noteinfo{以下说明依据本项目现有真题与模拟题整理，反映的是当前题库体现出的复习重点。}

\textbf{一、资料概况}

本项目当前收录本课程期中、期末题目各 5 套左右，覆盖往年真题与模拟题。现有资料显示，这门课以\textbf{偏微分方程的分离变量法、特殊函数与积分变换}为主线，题目通常兼顾“会列方程”和“会把展开系数真正算出来”。

\textbf{二、考试范围（按现有题库归纳）}
\begin{itemize}[leftmargin=2em,itemsep=0.2em,topsep=0.2em]
  \item 波动方程、热传导方程、Laplace 方程等典型方程的建模与定解问题；
  \item 直角、柱坐标、球坐标下的分离变量，以及本征值问题的建立；
  \item Sturm--Liouville 问题、本征函数正交性与模方计算；
  \item Fourier 级数 / Fourier 变换、Laplace 变换及其逆变换；
  \item Bessel 函数、Legendre 多项式 / 球谐函数等特殊函数在边值问题中的应用；
  \item 某些题目还涉及 Green 函数、卷积、积分表示或特殊恒等式的使用。
\end{itemize}

\textbf{三、内容与命题特点}

本课程常见命题方式是：\textbf{先让你判断用什么坐标与什么分离形式，再要求把本征值、边界条件和展开系数完整写出来}。因此它不是单纯的“套公式课”，而是很看重你对边界条件、正交基和坐标系选择的敏感度。

\textbf{四、难度评价}

整体难度可评为\textbf{中上到偏高}。难点主要不在计算量，而在于：
\begin{itemize}[leftmargin=2em,itemsep=0.2em,topsep=0.2em]
  \item 看清区域几何与边界条件后，能否迅速选对变量分离形式；
  \item 会不会把“本征值问题”和“原 PDE 的解”区分清楚；
  \item 在特殊函数出现时，是否知道哪些性质可以直接调用，哪些需要由边界条件推出。
\end{itemize}

\textbf{五、对课程的基本评价}

这门课非常锻炼结构化解题能力。题目虽然表面变化多，但本质上高度依赖少数标准套路：\textbf{建立方程、选坐标、分离变量、解本征问题、利用正交展开定系数}。只要把这条链条练熟，很多看起来很新的题都能迅速拆开。

\textbf{六、如何更好利用模拟题}
\begin{itemize}[leftmargin=2em,itemsep=0.2em,topsep=0.2em]
  \item 每做一题，先用一句话概括“为什么选这个坐标系”；
  \item 把所有本征值问题单独摘出来做一遍，训练看到边界条件就能判断正弦、余弦、Bessel 或球谐；
  \item 对计算较长的题，先练框架，再补系数细节；
  \item 做卷积、Laplace 变换类题时，整理常见变换对照表；
  \item 复习后期建议按“波动 / 热传导 / Laplace / 特殊函数 / 积分变换”五类回刷，而不是按年份刷。
\end{itemize}

